% =====================================================================
\chapter{Soluciones de los ejercicios}
\label{ap:soluciones}
% =====================================================================

Cada solución se presenta en cinco partes: datos, planteamiento, desarrollo, resultado y comprobación. En los pasos intermedios se redondea a cuatro cifras significativas salvo que se indique otra cosa; las potencias de~2 de los formatos numéricos son exactas. Los valores en coma flotante se han comprobado con PyTorch.

% ---------------------------------------------------------------------
\section{Tema 1}

\begin{solucion}{eje:1-congelada}
\datos red plana de 24 capas; factor por capa $1/\sqrt3$ con la inicialización por defecto (\cref{ej:pytorch}).
\planteamiento por la \cref{prop:grad-prof}, el gradiente de los pesos de la capa $\ell$ contiene el producto de los $L-\ell$ jacobianos que tiene encima.
\pasos
\begin{enumerate}[label=\arabic*.]
  \item El gradiente de la capa~1 atraviesa los jacobianos $J_2,\dots,J_{24}$: 23 factores.
  \item Si cada uno escala por $1/\sqrt3=0{,}5774$, el producto es $0{,}5774^{23}=3{,}3\cdot 10^{-6}$.
  \item El gradiente de la capa~24 no atraviesa ningún jacobiano: factor~1.
  \item Con la misma tasa de aprendizaje, el paso de la capa~1 es unas $3\cdot 10^{5}$ veces menor que el de la última: sus pesos apenas cambian.
\end{enumerate}
\resultado \textbf{la capa~1 recibe un gradiente del orden de $3\cdot 10^{-6}$ veces el de la última y queda prácticamente congelada}, aunque la última aprenda con normalidad.
\comprobacion la práctica del módulo mide un cociente de $1{,}4\cdot 10^{-5}$ entre las capas~1 y~24, del mismo orden de magnitud.
\end{solucion}

\begin{solucion}{eje:1-factores}
\datos 24 capas; factores $0{,}5$ y $1{,}1$; umbral $2^{-24}$.
\planteamiento el factor total es el producto de los factores por capa (\cref{ej:escalar}).
\pasos
\begin{enumerate}[label=\arabic*.]
  \item (a) $0{,}5^{24}=1/2^{24}=1/16\,777\,216=5{,}960\cdot 10^{-8}$.
  \item (b) $1{,}1^{2}=1{,}21$; $1{,}1^{4}=1{,}4641$; $1{,}1^{8}=2{,}1436$; $1{,}1^{16}=4{,}5950$; $1{,}1^{24}=1{,}1^{16}\cdot 1{,}1^{8}=4{,}5950\cdot 2{,}1436=9{,}850$.
  \item (c) $0{,}5^{L}<2^{-24}\iff 2^{-L}<2^{-24}\iff L>24$. Con $L=24$ el factor es exactamente $2^{-24}$; hace falta $L=25$.
\end{enumerate}
\resultado \textbf{(a) $5{,}96\cdot 10^{-8}$; (b) $9{,}85$; (c) 25 capas.}
\comprobacion por logaritmos, $1{,}1^{24}=e^{24\ln 1{,}1}=e^{24\cdot 0{,}09531}=e^{2{,}287}=9{,}85$.
\end{solucion}

\begin{solucion}{eje:1-saturacion}
\datos $\tanh'(t)=1-\tanh^{2}(t)$; 5 capas, $w_\ell=1$, preactivaciones~2; $s'(t)\le 1/4$.
\planteamiento el jacobiano escalar de cada capa es $f'(z_\ell)\,w_\ell$ (\cref{def:red}).
\pasos
\begin{enumerate}[label=\arabic*.]
  \item (a) $\tanh 0=0\Rightarrow\tanh'(0)=1$; $\tanh 1=0{,}7616$, $\tanh^{2}1=0{,}5800\Rightarrow\tanh'(1)=0{,}4200$.
  \item $\tanh 2=0{,}9640$, $\tanh^{2}2=0{,}9293\Rightarrow\tanh'(2)=0{,}07065$; $\tanh 3=0{,}9951$, $\tanh^{2}3=0{,}9901\Rightarrow\tanh'(3)=0{,}009866$.
  \item (b) Cada capa aporta $0{,}07065\cdot 1$; el producto es $0{,}07065^{5}=1{,}760\cdot 10^{-6}$.
  \item (c) La sigmoide alcanza su derivada máxima $1/4$ en $t=0$; en el mejor caso, $(1/4)^{5}=1/1024=9{,}766\cdot 10^{-4}$.
\end{enumerate}
\resultado \textbf{(a) $1$; $0{,}420$; $0{,}0707$; $0{,}0099$. (b) $1{,}76\cdot 10^{-6}$. (c) $9{,}77\cdot 10^{-4}$.}
\comprobacion el resultado (c) coincide con $\gamma^{L}$ del \cref{teo:cotas} con $\gamma=1/4$, $\norm{W}=1$ y $L=5$: incluso en el mejor caso, la sigmoide reduce el gradiente.
\end{solucion}

\begin{solucion}{eje:1-xavier}
\datos $n=64$; varianza inicial~1; $\sigma^{2}=1/(3n)$ por defecto.
\planteamiento \cref{cor:xavier} ($\sigma^{2}=1/n$), \cref{prop:he} ($\sigma^{2}=2/n$) y factor $n\sigma^{2}$ por capa.
\pasos
\begin{enumerate}[label=\arabic*.]
  \item (a) Red lineal: $\sigma^{2}=1/64=0{,}015625$, $\sigma=1/8=0{,}125$.
  \item Red ReLU: $\sigma^{2}=2/64=0{,}03125$, $\sigma=\sqrt{0{,}03125}=0{,}1768$.
  \item (b) $n\sigma^{2}=64\cdot 1/(3\cdot 64)=1/3$.
  \item Tras 10 capas: $(1/3)^{10}=1/59\,049=1{,}694\cdot 10^{-5}$.
\end{enumerate}
\resultado \textbf{(a) $\sigma=0{,}125$ (Xavier) y $\sigma\approx 0{,}177$ (He); (b) factor $1/3$ y varianza $1{,}69\cdot 10^{-5}$ tras 10 capas.}
\comprobacion He multiplica la varianza de Xavier por~2 y, por tanto, la desviación típica por $\sqrt2$: $0{,}125\cdot 1{,}414=0{,}1768$.
\end{solucion}

\begin{solucion}{eje:1-singulares}
\datos $A,B\in\R^{d\times d}$; $s_{\min}=0{,}9$, $s_{\max}=1{,}1$, $L=24$.
\planteamiento los valores singulares de una matriz $C$ son los valores $\norm{Cv}$ para sus vectores singulares derechos $v$ unitarios, de modo que todos están entre $\smin(C)$ y $\norm{C}_2$.
\pasos
\begin{enumerate}[label=\arabic*.]
  \item (a) Sea $\sigma$ un valor singular de $AB$ con vector singular derecho $v$, $\norm{v}=1$. Entonces $\sigma=\norm{ABv}$.
  \item $\norm{ABv}\le\norm{A}_2\norm{Bv}\le\norm{A}_2\norm{B}_2$, y $\norm{ABv}\ge\smin(A)\norm{Bv}\ge\smin(A)\smin(B)$ (\cref{lem:producto}).
  \item (b) Por inducción sobre el número de factores, $\smin(J_1\T\cdots J_L\T)\ge\prod_\ell\smin(J_\ell)\ge s_{\min}^{L}$ y $\norm{J_1\T\cdots J_L\T}_2\le s_{\max}^{L}$; el cociente del \cref{teo:cotas} queda entre ambos.
  \item $0{,}9^{24}=0{,}07977$ y $1{,}1^{24}=9{,}850$.
  \item (c) El intervalo $[0{,}080;\,9{,}85]$ es amplio aunque cada jacobiano esté a un 10\,\% de una isometría: el número de condición del producto puede llegar a $(1{,}1/0{,}9)^{24}=123{,}5$. Para un gradiente bien condicionado, todos los valores singulares de cada jacobiano deben estar muy próximos a~1.
\end{enumerate}
\resultado \textbf{los valores singulares de $AB$ están en $[\smin(A)\smin(B),\,\norm{A}_2\norm{B}_2]$; con 24 capas, el cociente está en $[0{,}080;\,9{,}85]$.}
\comprobacion con $A=\diag(2,1)$ y $B=\diag(1,3)$, $AB=\diag(2,3)$ tiene valores singulares $2$ y $3$, dentro de $[1\cdot 1,\,2\cdot 3]=[1,6]$.
\end{solucion}

\begin{solucion}{eje:1-potencias}
\datos $M=\sum_i\lambda_iv_iv_i\T$ con $v_i\T v_j=\delta_{ij}$; $M=\begin{psmallmatrix}0{,}85&0{,}35\\0{,}35&0{,}85\end{psmallmatrix}$, $x=(1,0)$.
\planteamiento descomposición espectral (módulo~3) y desarrollo de $x$ en la base de autovectores.
\pasos
\begin{enumerate}[label=\arabic*.]
  \item (a) Por inducción: si $M^{k}=\sum_i\lambda_i^{k}v_iv_i\T$, entonces $M^{k+1}=\sum_{i,j}\lambda_i^{k}\lambda_jv_i(v_i\T v_j)v_j\T=\sum_i\lambda_i^{k+1}v_iv_i\T$. Así $M^{k}v_i=\lambda_i^{k}v_i$.
  \item (b) Sea $x=\sum_ic_iv_i$ con $c_1\ne 0$. Entonces $M^{k}x=\lambda_1^{k}\bigl(c_1v_1+\sum_{i\ge 2}c_i(\lambda_i/\lambda_1)^{k}v_i\bigr)$.
  \item Como $\abs{\lambda_i/\lambda_1}<1$ para $i\ge 2$, los términos de la suma tienden a~0, y $M^{k}x/\norm{M^{k}x}\to\operatorname{signo}(c_1\lambda_1^{k})\,v_1$.
  \item (c) $x=(1,0)$ tiene $c_1=c_2=1/\sqrt2$ respecto de $v_1=(1,1)/\sqrt2$ y $v_2=(1,-1)/\sqrt2$. Para $k=1$, $Mx=(0{,}85;\,0{,}35)$, de dirección $(0{,}9247;\,0{,}3807)$.
  \item Para $k=10$, $M^{10}x=(3{,}0964;\,3{,}0954)$ (\cref{ej:potencias}), de norma $4{,}3782$ y dirección $(0{,}70722;\,0{,}70700)$, frente a $v_1=(0{,}70711;\,0{,}70711)$.
  \item (d) En una red recurrente lineal el gradiente contiene potencias $(W\T)^{k}$. Su componente sobre el autovector dominante escala como $\lambda_1^{k}$, que explota si $\abs{\lambda_1}>1$ y se desvanece si $\abs{\lambda_1}<1$; además, la dirección del gradiente queda dominada por $v_1$, con independencia de la señal de error.
\end{enumerate}
\resultado \textbf{$M^{k}$ tiene autovalores $\lambda_i^{k}$ y $M^{k}x$ se alinea con $v_1$; con $k=10$ la dirección coincide con $v_1$ en tres decimales.}
\comprobacion $1{,}2^{10}=6{,}1917$ y $\tfrac12(6{,}1917+0{,}00098)=3{,}0964$, $\tfrac12(6{,}1917-0{,}00098)=3{,}0954$.
\end{solucion}

% ---------------------------------------------------------------------
\section{Tema 2}

\begin{solucion}{eje:2-anidadas}
\planteamiento separar lo que se puede representar de lo que se puede aprender con datos y recursos finitos.
\pasos
\begin{enumerate}[label=\arabic*.]
  \item Razón estadística: una clase mayor contiene funciones que se ajustan a los datos de entrenamiento sin generalizar; con datos finitos, aumentar la capacidad puede aumentar el error de prueba (módulo~8).
  \item Razón computacional: el coste de cálculo y de memoria de cada paso crece linealmente con la profundidad, y la pasada hacia delante es secuencial a lo largo de las capas, lo que aumenta la latencia.
  \item Además, que la clase contenga una buena función no garantiza que el optimizador la encuentre; el flujo residual crece con la profundidad (\cref{prop:crecimiento}).
\end{enumerate}
\resultado \textbf{la inclusión de clases garantiza que la profundidad no empeora la capacidad de representación, pero no la generalización ni el coste.}
\end{solucion}

\begin{solucion}{eje:2-factor}
\datos 24 bloques; $J_\ell=\pm 0{,}01$; red plana con factor $0{,}5$.
\planteamiento en el caso escalar, el factor residual es $\prod_\ell(1+J_\ell)$ (\cref{teo:jac-residual}).
\pasos
\begin{enumerate}[label=\arabic*.]
  \item (a) $1{,}01^{2}=1{,}0201$; $1{,}01^{4}=1{,}0406$; $1{,}01^{8}=1{,}0829$; $1{,}01^{16}=1{,}1726$; $1{,}01^{24}=1{,}1726\cdot 1{,}0829=1{,}270$.
  \item (b) $0{,}99^{24}$: $0{,}99^{2}=0{,}9801$; $0{,}99^{4}=0{,}9606$; $0{,}99^{8}=0{,}9227$; $0{,}99^{16}=0{,}8515$; $0{,}99^{24}=0{,}8515\cdot 0{,}9227=0{,}7857$.
  \item (c) Red plana: $0{,}5^{24}=5{,}96\cdot 10^{-8}$. Red residual con $J_\ell=-0{,}5$: $(1-0{,}5)^{24}=5{,}96\cdot 10^{-8}$, idéntico.
\end{enumerate}
\resultado \textbf{(a) $1{,}27$; (b) $0{,}786$; (c) la red residual solo ayuda si las ramas son pequeñas; si todas se oponen a la identidad con $J_\ell=-0{,}5$, el factor es el mismo que en la red plana.}
\comprobacion $e^{\pm 0{,}24}=1{,}271$ y $0{,}787$, de acuerdo con $(1\pm x)^{n}\approx e^{\pm nx}$.
\end{solucion}

\begin{solucion}{eje:2-matrices}
\datos $J_1=\begin{psmallmatrix}0{,}2&0\\0{,}1&-0{,}1\end{psmallmatrix}$, $J_2=\begin{psmallmatrix}-0{,}1&0{,}3\\0&0{,}2\end{psmallmatrix}$.
\planteamiento \cref{teo:jac-residual} con $L=2$: $(I+J_2)(I+J_1)=I+J_1+J_2+J_2J_1$.
\pasos
\begin{enumerate}[label=\arabic*.]
  \item $I+J_2=\begin{psmallmatrix}0{,}9&0{,}3\\0&1{,}2\end{psmallmatrix}$ e $I+J_1=\begin{psmallmatrix}1{,}2&0\\0{,}1&0{,}9\end{psmallmatrix}$.
  \item Producto: fila~1, $(0{,}9\cdot 1{,}2+0{,}3\cdot 0{,}1;\ 0{,}3\cdot 0{,}9)=(1{,}11;\,0{,}27)$; fila~2, $(1{,}2\cdot 0{,}1;\ 1{,}2\cdot 0{,}9)=(0{,}12;\,1{,}08)$.
  \item $J_2J_1$: fila~1, $(-0{,}1\cdot 0{,}2+0{,}3\cdot 0{,}1;\ 0{,}3\cdot(-0{,}1))=(0{,}01;\,-0{,}03)$; fila~2, $(0{,}2\cdot 0{,}1;\ 0{,}2\cdot(-0{,}1))=(0{,}02;\,-0{,}02)$.
  \item $I+J_1+J_2=\begin{psmallmatrix}1{,}1&0{,}3\\0{,}1&1{,}1\end{psmallmatrix}$; sumando $J_2J_1$ se obtiene $\begin{psmallmatrix}1{,}11&0{,}27\\0{,}12&1{,}08\end{psmallmatrix}$.
\end{enumerate}
\resultado \textbf{$(I+J_2)(I+J_1)=\begin{psmallmatrix}1{,}11&0{,}27\\0{,}12&1{,}08\end{psmallmatrix}$, cercano a~$I$; la red plana daría $J_2J_1=\begin{psmallmatrix}0{,}01&-0{,}03\\0{,}02&-0{,}02\end{psmallmatrix}$}, con valores singulares $0{,}041$ y $0{,}0097$.
\comprobacion los valores singulares de la matriz residual son $1{,}293$ y $0{,}902$, mayores que $(1-a)^{2}=0{,}397$ con $a=\max(\norm{J_1}_2,\norm{J_2}_2)=\max(0{,}229;\,0{,}370)=0{,}370$, como exige la \cref{prop:cota-residual}.
\end{solucion}

\begin{solucion}{eje:2-flujo}
\datos $N=48$; $\Var(h_0)=1$; $v=0{,}5$; escalado $1/\sqrt N$.
\planteamiento \cref{prop:crecimiento}: $\Var(h_N)=\Var(h_0)+Nv$. Multiplicar unos pesos por $1/\sqrt N$ multiplica la varianza de su salida por $1/N$.
\pasos
\begin{enumerate}[label=\arabic*.]
  \item (a) $\Var(h_N)=1+48\cdot 0{,}5=1+24=25$; desviación típica $\sqrt{25}=5$.
  \item (b) Cada aportación tiene varianza $v/N=0{,}5/48=0{,}01042$.
  \item $\Var(h_N)=1+48\cdot 0{,}01042=1+0{,}5=1{,}5$; desviación típica $\sqrt{1{,}5}=1{,}225$.
\end{enumerate}
\resultado \textbf{(a) varianza~25, desviación típica~5; (b) varianza~$1{,}5$, desviación típica $1{,}22$.}
\comprobacion en (b) la varianza total añadida es $v=0{,}5$ con independencia de $N$: el escalado hace que la profundidad no cambie la escala del flujo.
\end{solucion}

\begin{solucion}{eje:2-cota}
\datos $h_\ell=h_{\ell-1}+F_\ell(h_{\ell-1})$; caso escalar con $\abs{J_\ell}\le a<1$.
\planteamiento desarrollo del producto de jacobianos y cota de cada factor.
\pasos
\begin{enumerate}[label=\arabic*.]
  \item (a) Cada bloque tiene jacobiano $I+J_\ell$; por la regla de la cadena, $\partial h_L/\partial h_0=(I+J_L)\cdots(I+J_1)$.
  \item Al desarrollar, cada término elige $I$ o $J_\ell$ en cada factor; el que elige siempre $I$ vale $I$. El desarrollo contiene, por tanto, el término $I$.
  \item (b) $\abs{J_\ell}\le a$ implica $J_\ell\ge -a$, luego $1+J_\ell\ge 1-a>0$.
  \item Todos los factores son positivos y mayores o iguales que $1-a$; multiplicando las $L$ desigualdades entre números positivos, $\prod_\ell(1+J_\ell)\ge(1-a)^{L}>0$. $\blacksquare$
  \item (c) Red plana con $J_1=0$ (unidad ReLU apagada) y $J_2=3$: $J_2J_1=0$. La red residual con las mismas ramas da $(1+3)(1+0)=4$.
\end{enumerate}
\resultado \textbf{el desarrollo contiene $I$; con $\abs{J_\ell}\le a<1$ el factor escalar es al menos $(1-a)^{L}>0$}; en una red plana basta una capa con jacobiano nulo para anularlo.
\comprobacion con $a=0{,}01$ y $L=24$, la cota es $0{,}99^{24}=0{,}786$, el valor del \cref{eje:2-factor}(b).
\end{solucion}

% ---------------------------------------------------------------------
\section{Tema 3}

\begin{solucion}{eje:3-clasificar}
\datos (a) $h+\operatorname{Attn}(\LN(h))$; (b) $\LN(h+\operatorname{MLP}(h))$.
\planteamiento \cref{def:pre-post}: pre-norm normaliza dentro de la rama; post-norm, después de la suma.
\pasos
\begin{enumerate}[label=\arabic*.]
  \item (a) La normalización está dentro de la rama: \textbf{pre-norm}. Su jacobiano es $I+J_{\mathrm{Attn}}N$, con el término $I$.
  \item (b) La normalización envuelve la suma: \textbf{post-norm}. Su jacobiano es $N(I+J_{\mathrm{MLP}})$; la identidad queda multiplicada por $N$.
  \item Con $L$ bloques post-norm, el camino directo atraviesa $L$ jacobianos de normalización y deja de ser la identidad (\cref{prop:jac-pre-post}).
\end{enumerate}
\resultado \textbf{(a) pre-norm; (b) post-norm. Solo (a) conserva la identidad en el camino residual.}
\comprobacion es el patrón de la \cref{fig:prenorm}.
\end{solucion}

\begin{solucion}{eje:3-preactivacion}
\planteamiento identificar qué operaciones quedan en el camino principal, fuera de la rama.
\pasos
\begin{enumerate}[label=\arabic*.]
  \item Ordenación original: $h_\ell=\ReLU\bigl(h_{\ell-1}+F(h_{\ell-1})\bigr)$. La ReLU posterior a la suma está en el camino principal; su jacobiano $\diag(\mathbf 1[\,\cdot>0])$ multiplica la identidad y puede anular componentes.
  \item Ordenación de preactivación: $h_\ell=h_{\ell-1}+F(h_{\ell-1})$ con $F=\mathrm{conv}\circ\ReLU\circ\mathrm{BN}$. Tras la suma no hay nada: el camino principal es exactamente la identidad y se aplica el \cref{teo:jac-residual}.
  \item Pre-norm hace lo mismo con la normalización de un Transformer: la mueve dentro de la rama y deja el flujo residual libre.
\end{enumerate}
\resultado \textbf{la preactivación saca del camino principal la normalización y la no linealidad, con lo que el jacobiano de cada bloque es $I+J_\ell$; es la misma idea que pre-norm.}
\end{solucion}

\begin{solucion}{eje:3-calculo}
\datos $x=(1,2,3,6)$, $d=4$, $\gamma=g=\uno$, $\beta=0$, $\epsilon=0$.
\planteamiento \cref{def:layernorm,def:rmsnorm}.
\pasos
\begin{enumerate}[label=\arabic*.]
  \item Media: $\mu=(1+2+3+6)/4=12/4=3$.
  \item Desviaciones: $x-\mu=(-2,-1,0,3)$; cuadrados $(4,1,0,9)$; $\sigma^{2}=14/4=3{,}5$; $\sigma=1{,}871$.
  \item $\LN(x)=(-2/1{,}871;\ -1/1{,}871;\ 0;\ 3/1{,}871)=(-1{,}069;\,-0{,}5345;\,0;\,1{,}604)$.
  \item Cuadrados de $x$: $(1,4,9,36)$; media $50/4=12{,}5$; $\RMS(x)=\sqrt{12{,}5}=3{,}536$.
  \item $\RMSNorm(x)=(0{,}2828;\,0{,}5657;\,0{,}8485;\,1{,}697)$.
\end{enumerate}
\resultado \textbf{$\LN(x)\approx(-1{,}069;\,-0{,}5345;\,0;\,1{,}604)$;}\newline \textbf{$\RMSNorm(x)\approx(0{,}2828;\,0{,}5657;\,0{,}8485;\,1{,}697)$.}
\comprobacion la salida de LayerNorm suma $-1{,}069-0{,}5345+0+1{,}604\approx 0$ y la media de sus cuadrados es $(1{,}143+0{,}2857+0+2{,}571)/4=1{,}000$. La media de los cuadrados de RMSNorm es $(0{,}08+0{,}32+0{,}72+2{,}88)/4=1$.
\end{solucion}

\begin{solucion}{eje:3-invariancia}
\datos $x=(1,2,3,6)$; $2x+10\uno=(12,14,16,22)$; $x+3\uno=(4,5,6,9)$.
\planteamiento recalcular los estadísticos y comparar con el \cref{eje:3-calculo}.
\pasos
\begin{enumerate}[label=\arabic*.]
  \item (a) $\mu=64/4=16$; desviaciones $(-4,-2,0,6)$; cuadrados $(16,4,0,36)$; $\sigma^{2}=56/4=14$; $\sigma=3{,}742$.
  \item $\LN(2x+10\uno)=(-4,-2,0,6)/3{,}742=(-1{,}069;\,-0{,}5345;\,0;\,1{,}604)$, igual que $\LN(x)$.
  \item (b) Cuadrados de $x+3\uno$: $(16,25,36,81)$; media $158/4=39{,}5$; $\RMS=6{,}285$.
  \item $\RMSNorm(x+3\uno)=(0{,}6364;\,0{,}7956;\,0{,}9547;\,1{,}432)$, distinto de $\RMSNorm(x)=(0{,}2828;\,0{,}5657;\,0{,}8485;\,1{,}697)$.
\end{enumerate}
\resultado \textbf{(a) coinciden: invariancia de LayerNorm frente a escala y traslación (\cref{prop:invariancias}(i)); (b) no coinciden: RMSNorm no es invariante frente a traslaciones (\cref{prop:invariancias}(ii)).}
\comprobacion en (a) la desviación típica se multiplica por $\alpha=2$: $3{,}742=2\cdot 1{,}871$.
\end{solucion}

\begin{solucion}{eje:3-jacobiano}
\datos $x=(3,4)$, $d=2$, $g=\uno$, $\epsilon=0$.
\planteamiento \cref{teo:jac-rms}: $J=\frac1r\bigl(I-xx\T/\norm{x}^{2}\bigr)$.
\pasos
\begin{enumerate}[label=\arabic*.]
  \item $\norm{x}^{2}=9+16=25$; $r=\sqrt{25/2}=3{,}536$.
  \item $xx\T/25=\begin{psmallmatrix}9&12\\12&16\end{psmallmatrix}/25=\begin{psmallmatrix}0{,}36&0{,}48\\0{,}48&0{,}64\end{psmallmatrix}$; $I-xx\T/25=\begin{psmallmatrix}0{,}64&-0{,}48\\-0{,}48&0{,}36\end{psmallmatrix}$.
  \item $J=\frac{1}{3{,}536}\begin{psmallmatrix}0{,}64&-0{,}48\\-0{,}48&0{,}36\end{psmallmatrix}=\begin{psmallmatrix}0{,}1810&-0{,}1358\\-0{,}1358&0{,}1018\end{psmallmatrix}$.
  \item $Jx=\frac1r(0{,}64\cdot 3-0{,}48\cdot 4;\ -0{,}48\cdot 3+0{,}36\cdot 4)=\frac1r(1{,}92-1{,}92;\ -1{,}44+1{,}44)=(0,0)$.
\end{enumerate}
\resultado \textbf{$J\approx\begin{psmallmatrix}0{,}1810&-0{,}1358\\-0{,}1358&0{,}1018\end{psmallmatrix}$, simétrica y con $Jx=0$.}
\comprobacion la diferenciación automática de PyTorch da la misma matriz. Además, $(-4,3)\perp x$ y $J(-4,3)=(-4,3)/r$: en la dirección ortogonal, $J$ actúa como $1/r$ por la identidad.
\end{solucion}

\begin{solucion}{eje:3-ortogonal}
\datos $y=\RMSNorm(x)$, $\epsilon=0$ salvo en (d); $\Loss=\Loss(y)$.
\planteamiento \cref{teo:jac-rms} y regla de la cadena.
\pasos
\begin{enumerate}[label=\arabic*.]
  \item (a) $\RMS(\alpha x)=\sqrt{\alpha^{2}\frac1d\sum x_i^{2}}=\alpha\RMS(x)$ para $\alpha>0$, luego $g_i\alpha x_i/(\alpha\RMS(x))=\RMSNorm(x)_i$.
  \item Derivando $\phi(\alpha)=\RMSNorm(\alpha x)$, constante, en $\alpha=1$: $0=\phi'(1)=J\,x$.
  \item (b) $\nabla_x\Loss=J\T\nabla_y\Loss$ y $x\T J\T\nabla_y\Loss=(Jx)\T\nabla_y\Loss=0$.
  \item (c) $\frac{d}{dt}\norm{x}^{2}=2x\T\dot x=-2x\T\nabla_x\Loss=0$: la norma es constante.
  \item Con un paso discreto, $\norm{x-\eta\nabla_x\Loss}^{2}=\norm{x}^{2}-2\eta\,x\T\nabla_x\Loss+\eta^{2}\norm{\nabla_x\Loss}^{2}=\norm{x}^{2}+\eta^{2}\norm{\nabla_x\Loss}^{2}$, mayor que $\norm{x}^{2}$ si el gradiente no es nulo.
  \item (d) Con $\epsilon>0$, $r^{2}=s+\epsilon$ con $s=\norm{x}^{2}/d$. Repitiendo el cálculo, $(Jx)_i=\frac{g_ix_i}{r}\bigl(1-\frac{s}{r^{2}}\bigr)=\frac{g_ix_i\,\epsilon}{r^{3}}$, que no es nulo pero es despreciable si $s\gg\epsilon$.
\end{enumerate}
\resultado \textbf{$Jx=0$ y $x\T\nabla_x\Loss=0$; la norma se conserva en el flujo continuo y crece en cada paso discreto. Con $\epsilon>0$ la invariancia es solo aproximada.}
\comprobacion en el \cref{eje:3-jacobiano}, $x\T J\T v=(Jx)\T v=0$ para todo $v$, porque $Jx=0$ y $J$ es simétrica.
\end{solucion}

% ---------------------------------------------------------------------
\section{Tema 4}

\begin{solucion}{eje:4-eval}
\planteamiento el dropout invertido conserva la esperanza, pero cada máscara es aleatoria (\cref{prop:dropout}).
\pasos
\begin{enumerate}[label=\arabic*.]
  \item Falta poner el modelo en modo evaluación (\texttt{model.eval()} en PyTorch), que desactiva el dropout.
  \item Con dropout activo, cada pasada sortea una máscara distinta y la salida cambia.
  \item La media no lo delata: $\E[\tilde h_i]=(1-p)\,h_i/(1-p)+p\cdot 0=h_i$.
  \item Lo que cambia es la varianza, $\frac{p}{1-p}h_i^{2}$: cada ejecución concreta tiene ruido.
\end{enumerate}
\resultado falta \texttt{model.eval()}; \textbf{el dropout invertido es insesgado, por lo que la media coincide, pero cada ejecución tiene ruido.}
\comprobacion con $p=0{,}3$ y $h_i=5$, $0{,}7\cdot 7{,}143+0{,}3\cdot 0=5{,}000$ (\cref{eje:4-mascara}).
\end{solucion}

\begin{solucion}{eje:4-constante}
\datos $P(d_i=0)=p$; $p=0{,}2$; $h=(1,3,0,2)$; $d=(1,0,1,1)$.
\planteamiento $\E[d_i]=1-p$ y linealidad de la esperanza.
\pasos
\begin{enumerate}[label=\arabic*.]
  \item (a) $\E[h_{\mathrm{drop},i}]=\gamma h_i\E[d_i]=\gamma h_i(1-p)$. Igualando a $h_i$ para $h_i\ne 0$: $\gamma=1/(1-p)$.
  \item (b) $\gamma=1/0{,}8=1{,}25$; $d\odot h=(1,0,0,2)$; $h_{\mathrm{drop}}=(1{,}25;\,0;\,0;\,2{,}5)$.
  \item (c) Al entrenar, la máscara impide que la red dependa de unidades concretas y actúa como regularizador. Al evaluar se quiere una predicción determinista, la del modelo medio; gracias a $\gamma$, usar $h$ sin máscara da la esperanza correcta.
\end{enumerate}
\resultado \textbf{$\gamma=1/(1-p)=1{,}25$; $h_{\mathrm{drop}}=(1{,}25;\,0;\,0;\,2{,}5)$.} Es el dropout invertido de la \cref{def:dropout}.
\comprobacion $\E[h_{\mathrm{drop},2}]=0{,}8\cdot 1{,}25\cdot 3+0{,}2\cdot 0=3=h_2$.
\end{solucion}

\begin{solucion}{eje:4-mascara}
\datos $p=0{,}3$; $h=(5,5,5,5)$; $m=(1,0,1,1)$.
\planteamiento $\tilde h_i=m_ih_i/(1-p)$.
\pasos
\begin{enumerate}[label=\arabic*.]
  \item $1-p=0{,}7$; factor $1/0{,}7=1{,}429$.
  \item $m\odot h=(5,0,5,5)$; $\tilde h=(7{,}143;\,0;\,7{,}143;\,7{,}143)$.
  \item $\E[\tilde h_i]=0{,}7\cdot(5/0{,}7)+0{,}3\cdot 0=5=h_i$.
  \item Media de la salida para esta máscara: $3\cdot 7{,}143/4=5{,}357$.
\end{enumerate}
\resultado \textbf{$\tilde h\approx(7{,}143;\,0;\,7{,}143;\,7{,}143)$, con $\E[\tilde h_i]=5$; la media de esta realización es $5{,}357$.}
\comprobacion con una fracción $3/4$ de supervivientes, la media es $5\cdot(3/4)/0{,}7=5{,}357$; promediando sobre muchas máscaras la fracción tiende a $0{,}7$ y la media a~5.
\end{solucion}

\begin{solucion}{eje:4-recorte}
\datos $g=(6,8)$; $c=5$ y $c=20$.
\planteamiento $\hat g=g\min(1,c/\norm{g})$ (\cref{def:recorte}).
\pasos
\begin{enumerate}[label=\arabic*.]
  \item $\norm{g}=\sqrt{36+64}=\sqrt{100}=10$.
  \item $c=5$: $\alpha=\min(1;\,0{,}5)=0{,}5$; $\hat g=(3,4)$.
  \item $c=20$: $\alpha=\min(1;\,2)=1$; $\hat g=(6,8)$.
\end{enumerate}
\resultado \textbf{$c=5$: $\hat g=(3,4)$; $c=20$: $\hat g=(6,8)$, sin recorte.}
\comprobacion $\norm{(3,4)}=5=\min(10,5)$ y la dirección $(0{,}6;\,0{,}8)$ se conserva, como afirma la \cref{prop:recorte}(i).
\end{solucion}

\begin{solucion}{eje:4-global}
\datos $g^{(1)}=(1,2)$, $g^{(2)}=(2)$, $c=1$.
\planteamiento la norma global es la del vector concatenado; el mismo factor se aplica a todos los tensores.
\pasos
\begin{enumerate}[label=\arabic*.]
  \item (a) $g=(1,2,2)$, $\norm{g}=\sqrt{1+4+4}=3$; $\alpha=1/3$.
  \item $\hat g^{(1)}=(0{,}3333;\,0{,}6667)$ y $\hat g^{(2)}=(0{,}6667)$.
  \item (b) $\norm{g^{(1)}}=\sqrt5=2{,}236$, luego $\hat g^{(1)}=(0{,}4472;\,0{,}8944)$; $\norm{g^{(2)}}=2$, luego $\hat g^{(2)}=(1)$.
  \item El vector $(0{,}4472;\,0{,}8944;\,1)$ no es proporcional a $(1,2,2)$: el cociente entre la tercera y la primera componente es $2{,}236$ en lugar de~2.
\end{enumerate}
\resultado \textbf{norma global: $\hat g^{(1)}=(1/3,\,2/3)$, $\hat g^{(2)}=(2/3)$. Recortar por separado cambia la dirección del gradiente total.}
\comprobacion $\norm{(1/3,2/3,2/3)}=\sqrt{1/9+4/9+4/9}=1=c$.
\end{solucion}

\begin{solucion}{eje:4-epsilon}
\datos $g^{(1)}=(2,-1)$, $g^{(2)}=(2)$, $M=1$, $\epsilon=10^{-6}$.
\planteamiento factor $M/(\norm{g}+\epsilon)$ si $\norm{g}\ge M$.
\pasos
\begin{enumerate}[label=\arabic*.]
  \item (a) $\norm{g}=\sqrt{4+1+4}=3\ge 1$: hay que recortar.
  \item Factor: $1/(3+10^{-6})=0{,}3333332$.
  \item $\hat g^{(1)}=(0{,}6667;\,-0{,}3333)$ y $\hat g^{(2)}=(0{,}6667)$.
  \item (b) Norma resultante: $3/(3+10^{-6})=0{,}9999997$.
  \item El $\epsilon$ evita la división por cero cuando $g=0$ y deja la norma ligeramente por debajo de $M$.
\end{enumerate}
\resultado \textbf{norma global~3 y factor $\approx 1/3$; $\hat g^{(1)}\approx(0{,}6667;\,-0{,}3333)$, $\hat g^{(2)}\approx(0{,}6667)$; norma final $0{,}9999997$.}
\comprobacion \texttt{torch.nn.utils.clip\_grad\_norm\_} devuelve la norma original, 3, y deja los mismos gradientes.
\end{solucion}

\begin{solucion}{eje:4-varianza}
\datos $w=(1,2,-1)$, $h=(2,1,2)$, $p=0{,}5$; $m_i\sim\Bern(0{,}5)$ independientes.
\planteamiento \cref{prop:dropout} y \cref{cor:dropout-lineal}; enumeración de las $2^{3}$ máscaras equiprobables.
\pasos
\begin{enumerate}[label=\arabic*.]
  \item (a) Las demostraciones son las del texto: $\E[m_i]=\E[m_i^{2}]=1-p$ da $\E[\tilde h_i]=h_i$ y $\E[\tilde h_i^{2}]=h_i^{2}/(1-p)$, de donde $\Var(\tilde h_i)=h_i^{2}p/(1-p)$; la independencia de las máscaras da la suma de varianzas.
  \item (b) $\E[y]=w\T h=1\cdot 2+2\cdot 1+(-1)\cdot 2=2$.
  \item $\Var(y)=\frac{0{,}5}{0{,}5}(1\cdot 4+4\cdot 1+1\cdot 4)=12$.
  \item Enumeración: $w_ih_i/(1-p)=(4,4,-4)$, luego $y=4(m_1+m_2-m_3)$. Las ocho máscaras dan $0,\,4,\,4,\,-4,\,8,\,0,\,0,\,4$, cada una con probabilidad $1/8$.
  \item Media: $16/8=2$. Segundo momento: $(0+16+16+16+64+0+0+16)/8=128/8=16$. Varianza: $16-2^{2}=12$.
  \item (c) Sin dropout, la activación $h_i$ es fija dada la entrada; con dropout, su segundo momento pasa a $h_i^{2}/(1-p)$ y su varianza a $\frac{p}{1-p}h_i^{2}$, que tiende a infinito cuando $p\to 1$ porque las pocas unidades supervivientes se amplifican por $1/(1-p)$.
\end{enumerate}
\resultado \textbf{$\E[y]=2$ y $\Var(y)=12$, confirmados por enumeración; la varianza que añade el dropout es $\frac{p}{1-p}h_i^{2}$ por componente.}
\comprobacion la desviación típica es $\sqrt{12}=3{,}464$, mayor que la media: con $p=0{,}5$ el ruido es grande.
\end{solucion}

% ---------------------------------------------------------------------
\section{Tema 5}

\begin{solucion}{eje:5-riesgo}
\datos $70\,000$ y $1{,}003$; fp16 ($x_{\max}=65\,504$, $\varepsilon=2^{-10}$), bf16 ($x_{\max}=3{,}39\cdot 10^{38}$, $\varepsilon=2^{-7}$), fp32.
\planteamiento desbordamiento si el valor supera $x_{\max}$; pérdida de precisión según la separación cerca de~1, que es $\varepsilon$.
\pasos
\begin{enumerate}[label=\arabic*.]
  \item $70\,000>65\,504$: en fp16 se convierte en \Inf. En bf16 cabe, redondeado a $70\,144$ (separación $2^{16}\cdot 2^{-7}=512$). En fp32 es exacto.
  \item Cerca de~1, bf16 separa los números $2^{-7}=0{,}0078$; los vecinos de $1{,}003$ son $1$ y $1{,}0078$, y el más cercano es~1. Error: $0{,}003$.
  \item fp16 los separa $2^{-10}=0{,}000977$; $1{,}003$ se redondea a $1+3\cdot 2^{-10}=1{,}00293$. Error: $7\cdot 10^{-5}$.
\end{enumerate}
\resultado \textbf{$70\,000$ es arriesgado en fp16 (da \Inf); $1{,}003$ pierde más precisión en bf16 (se almacena como~1).}
\comprobacion en PyTorch, $1{,}003$ en bf16 da $1{,}0$ y en fp16 da $1{,}0029$.
\end{solucion}

\begin{solucion}{eje:5-deducir}
\datos fp16: $e=5$, $m=10$; bf16: $e=8$, $m=7$.
\planteamiento \cref{prop:rango}: $b=2^{e-1}-1$, $x_{\max}=(2-2^{-m})2^{b}$, $x_{\min}=2^{1-b}$, $\varepsilon=2^{-m}$.
\pasos
\begin{enumerate}[label=\arabic*.]
  \item fp16: $b=2^{4}-1=15$.
  \item $x_{\max}=(2-2^{-10})2^{15}=2^{16}-2^{5}=65\,536-32=65\,504$.
  \item $x_{\min}=2^{-14}=1/16\,384=6{,}104\cdot 10^{-5}$; $\varepsilon=2^{-10}=9{,}766\cdot 10^{-4}$.
  \item bf16: $b=2^{7}-1=127$.
  \item $x_{\max}=(2-2^{-7})2^{127}=1{,}9922\cdot 1{,}7014\cdot 10^{38}=3{,}390\cdot 10^{38}$.
  \item $x_{\min}=2^{-126}=1{,}175\cdot 10^{-38}$; $\varepsilon=2^{-7}=0{,}0078125$.
\end{enumerate}
\resultado \textbf{fp16: $65\,504$, $6{,}10\cdot 10^{-5}$, $9{,}77\cdot 10^{-4}$. bf16: $3{,}39\cdot 10^{38}$, $1{,}18\cdot 10^{-38}$, $0{,}0078125$.}
\comprobacion bf16 tiene casi el rango de fp32 ($3{,}40\cdot 10^{38}$ y $1{,}18\cdot 10^{-38}$) porque comparten los 8 bits de exponente.
\end{solucion}

\begin{solucion}{eje:5-fp8}
\datos E5M2: $e=5$, $m=2$; E4M3: $e=4$, $m=3$.
\planteamiento \cref{prop:rango} y \cref{obs:e4m3}.
\pasos
\begin{enumerate}[label=\arabic*.]
  \item (a) E5M2: $b=15$; $x_{\max}=(2-2^{-2})2^{15}=1{,}75\cdot 32\,768=57\,344$; $x_{\min}=2^{-14}=6{,}10\cdot 10^{-5}$; $\varepsilon=2^{-2}=0{,}25$.
  \item (b) E4M3 con reglas IEEE: $b=2^{3}-1=7$; el mayor exponente normal sería $E=14$, con $E-b=7$; $x_{\max}=(2-2^{-3})2^{7}=1{,}875\cdot 128=240$; $x_{\min}=2^{-6}=0{,}015625$; $\varepsilon=2^{-3}=0{,}125$.
  \item E4M3 real: sin infinitos y con un único patrón \NaN\ (todos los bits de exponente y fracción a~1), el exponente $E=15$ es normal salvo en ese patrón. La mayor mantisa permitida con $E=15$ es $(1{,}110)_2=1{,}75$, y $x_{\max}=1{,}75\cdot 2^{8}=448$.
  \item (c) Los gradientes abarcan muchos órdenes de magnitud y requieren rango: E5M2. Pesos y activaciones requieren más precisión relativa: E4M3.
\end{enumerate}
\resultado \textbf{E5M2: $57\,344$, $6{,}10\cdot 10^{-5}$, $0{,}25$. E4M3: con reglas IEEE serían $240$, $0{,}0156$, $0{,}125$; el máximo real es~448.}
\comprobacion \texttt{torch.finfo} da $448$ y $57\,344$ como máximos de \texttt{float8\_e4m3fn} y \texttt{float8\_e5m2}.
\end{solucion}

\begin{solucion}{eje:5-mil}
\datos bf16 en $[512,1024)$: separación~4; fp16 en $[512,1024)$: separación $0{,}5$.
\planteamiento el resultado exacto se redondea al representable más cercano; en caso de empate, al de mantisa par (\cref{def:fl}).
\pasos
\begin{enumerate}[label=\arabic*.]
  \item (a) $1000+1=1001$; vecinos representables en bf16: $1000$ y $1004$, a distancias~1 y~3. Resultado: $1000$.
  \item $1000+4=1004=4\cdot 251$, representable: resultado exacto.
  \item (b) En fp16, $1001=0{,}5\cdot 2002$ es representable: resultado exacto.
  \item (c) Por el \cref{lem:absorcion}, se pierde todo sumando $0\le y<2$, la mitad de la separación. Con $y=2$ hay empate entre $1000=4\cdot 250$ y $1004=4\cdot 251$, y gana el de mantisa par, $1000$; con $y>2$ el resultado es $1004$.
\end{enumerate}
\resultado \textbf{bf16: $1000+1=1000$ y $1000+4=1004$; fp16: $1000+1=1001$. En bf16 se pierde cualquier sumando de hasta~2.}
\comprobacion PyTorch da $1000$ para $1000+2$ y $1004$ para $1000+2{,}5$ en bf16.
\end{solucion}

\begin{solucion}{eje:5-rmsnorm}
\datos $x=(300,400)$ en fp16 ($x_{\max}=65\,504$); $g=\uno$; $\epsilon=10^{-5}$; $d=2$.
\planteamiento $\RMSNorm(x)_i=x_i/\sqrt{\frac1d\sum x_j^{2}+\epsilon}$; el cuadrado es lo que desborda.
\pasos
\begin{enumerate}[label=\arabic*.]
  \item (a) En fp16: $300^{2}=90\,000>65\,504\to\Inf$; $400^{2}=160\,000\to\Inf$.
  \item Media $\Inf$; $\RMS=\sqrt{\Inf+10^{-5}}=\Inf$; salida $(300/\Inf,\ 400/\Inf)=(0,0)$. La señal se destruye sin ningún error.
  \item (b) En fp32: $x^{2}=(90\,000;\ 160\,000)$; media $125\,000$; $\RMS=\sqrt{125\,000{,}00001}=353{,}6$.
  \item Salida: $(300/353{,}6;\ 400/353{,}6)=(0{,}8485;\,1{,}131)$.
  \item En fp16, con separación $2^{-11}$ en $[0{,}5;1)$ y $2^{-10}$ en $[1,2)$: $(0{,}8486;\,1{,}132)$.
\end{enumerate}
\resultado \textbf{todo en fp16: $(0,0)$; en fp32 convertido a fp16: $(0{,}8486;\,1{,}132)$.}
\comprobacion $300/\sqrt{125\,000}=3/(5\sqrt2)=0{,}8485$ y $4/(5\sqrt2)=1{,}131$; la media de sus cuadrados es $(0{,}72+1{,}28)/2=1$.
\end{solucion}

\begin{solucion}{eje:5-absorcion}
\datos formato con $m$ bits de mantisa; fp16, $m=10$.
\planteamiento \cref{prop:rango}(iii): en $[2^{k},2^{k+1})$ la separación es $2^{k-m}$.
\pasos
\begin{enumerate}[label=\arabic*.]
  \item (a) Sea $2^{k}\le x<2^{k+1}$. Los representables vecinos de $x$ distan $2^{k-m}$ entre sí, luego $\abs{\fl(x)-x}\le 2^{k-m-1}$. Dividiendo por $x\ge 2^{k}$: $\abs\delta\le 2^{-m-1}=u$.
  \item (b) El representable siguiente a $x$ es $x+2^{k-m}$. Si $0\le y<2^{k-m-1}$, entonces $\abs{(x+y)-x}=y<2^{k-m}-y=\abs{(x+2^{k-m})-(x+y)}$, y $\fl(x+y)=x$.
  \item (c) Sea $w\in[2^{k},2^{k+1})$. Si $w/\abs{\Delta w}>4096=2^{12}$, entonces $\abs{\Delta w}<w/2^{12}<2^{k+1}/2^{12}=2^{k-11}$, la mitad de la separación en fp16, y por (b) la actualización se pierde.
  \item Ejemplo: $w=1$ ($k=0$), semiseparación $2^{-11}=4{,}883\cdot 10^{-4}$. Con $\Delta w=4{,}8\cdot 10^{-4}$, $w/\Delta w=2083>2048$ y $\fl(1+0{,}00048)=1$.
\end{enumerate}
\resultado \textbf{$\abs\delta\le u$; un sumando menor que la semiseparación se absorbe; en fp16 una actualización se pierde siempre si $w/\abs{\Delta w}>4096$, y puede perderse desde $2048$.} $\blacksquare$
\comprobacion en PyTorch, $1+0{,}00048$ en fp16 da $1$, mientras que $1+0{,}0005$ da $1{,}000977$, porque $0{,}0005$ supera la semiseparación.
\end{solucion}

\begin{solucion}{eje:5-acumular}
\datos mil sumas de $0{,}01$; fp16 con $m=10$; en $[2^{k},2^{k+1})$ la separación es $2^{k-10}$.
\planteamiento cada suma se redondea al múltiplo más próximo de la separación del intervalo en que está $s$; el error por suma depende de cuántas separaciones mide el sumando.
\pasos
\begin{enumerate}[label=\arabic*.]
  \item (a) $0{,}01\in[2^{-7},2^{-6})$, con separación $2^{-7-10}=2^{-17}=7{,}629\cdot 10^{-6}$. Como $0{,}01/2^{-17}=1310{,}72$, se redondea a $1311$ separaciones: $1311\cdot 2^{-17}=0{,}0100021$.
  \item (b) $s\in[4,8)$: separación $2^{-8}=0{,}003906$; el sumando mide $2{,}561$ separaciones y se redondea a~3, de modo que cada suma aumenta $s$ en $3\cdot 2^{-8}=0{,}01172$.
  \item $s\in[8,16)$: separación $2^{-7}=0{,}0078125$; el sumando mide $1{,}280$ separaciones y se redondea a~1: cada suma aumenta $s$ en $0{,}0078125$.
  \item (c) (i) En fp32 la separación cerca de~10 es $2^{3-23}=9{,}5\cdot 10^{-7}$ y el error acumulado es pequeño: $10{,}0001$.
  \item (ii) En fp16 el error por suma cambia de signo según el intervalo. En $[0{,}5;4)$ cada suma aumenta $s$ en $0{,}009766$, por defecto; en $[4,8)$, en $0{,}01172$, por exceso; en $[8,10)$, en $0{,}0078125$, por defecto. La simulación paso a paso da 359 sumas en $[0{,}5;4)$ con error total $-0{,}0830$; 341 en $[4,8)$ con $+0{,}5815$; 250 en $[8,10)$ con $-0{,}5474$; y 50 por debajo de $0{,}5$ con $-0{,}0001$.
  \item El neto es $-0{,}0490$, y $1000\cdot 0{,}0100021-0{,}0490=10{,}0021-0{,}0490=9{,}953$.
  \item (iii) y (iv): el acumulador fp32 apenas pierde nada, pero el sumando ya es el $0{,}01$ de fp16: $1000\cdot 0{,}0100021=10{,}0021$. En (iii) la conversión a fp32 se hace implícitamente antes de sumar; en (iv), explícitamente.
\end{enumerate}
\resultado \textbf{(a) $0{,}01\approx 1311\cdot 2^{-17}$; (b) $+0{,}01172$ por suma en $[4,8)$ y $+0{,}0078125$ en $[8,16)$; (c) el error de (ii) procede del acumulador fp16; el de (iii) y (iv), solo del redondeo del sumando.}
\comprobacion $359+341+250+50=1000$ sumas. Es el argumento de la acumulación en fp32 de la precisión mixta y del \cref{teo:suma}.
\end{solucion}

% ---------------------------------------------------------------------
\section{Tema 6}

\begin{solucion}{eje:6-diagnostico}
\planteamiento \cref{tab:diagnostico} y \cref{fig:diagnostico}.
\pasos
\begin{enumerate}[label=(\alph*)]
  \item Salto brusco a \Inf\ en fp16 sin crecimiento previo: desbordamiento de alguna activación o gradiente por encima de $65\,504$. Remedio: bf16, o escalado dinámico de la pérdida que omita el paso y reduzca $S$.
  \item Crecimiento de varios órdenes y después \NaN: gradientes que explotan. Remedio: bajar la tasa de aprendizaje, alargar el calentamiento y recortar el gradiente.
  \item \texttt{torch.log(torch.softmax(z))}: alguna probabilidad se anula por subdesbordamiento y $\ln 0=-\infty$; en la retropropagación aparecen operaciones como $0\cdot\infty$ que dan \NaN. Remedio: \texttt{log\_softmax} o la entropía cruzada aplicada a los logits (\cref{prop:ce}).
\end{enumerate}
\resultado \textbf{(a) desbordamiento en fp16; (b) gradientes que explotan; (c) logaritmo de una probabilidad nula.}
\end{solucion}

\begin{solucion}{eje:6-lse}
\datos $z=(1000,999,998)$.
\planteamiento \cref{prop:lse-estable} con $c=1000$.
\pasos
\begin{enumerate}[label=\arabic*.]
  \item $z-c\uno=(0,-1,-2)$.
  \item $e^{0}=1$; $e^{-1}=0{,}3679$; $e^{-2}=0{,}1353$; $S=1{,}503$.
  \item $\ln 1{,}503=0{,}4076$; $\LSE(z)=1000{,}408$.
  \item $\softmax(z)=(1/1{,}503;\ 0{,}3679/1{,}503;\ 0{,}1353/1{,}503)=(0{,}665;\,0{,}245;\,0{,}090)$.
\end{enumerate}
\resultado \textbf{$\LSE(z)\approx 1000{,}408$ y $\softmax(z)\approx(0{,}665;\,0{,}245;\,0{,}090)$.}
\comprobacion las probabilidades suman $1{,}000$ y $1000\le 1000{,}408\le 1000+\ln 3=1001{,}099$.
\end{solucion}

\begin{solucion}{eje:6-softmax}
\datos $v=(100,101,103)$ en fp32; $\ln x_{\max}=88{,}72$.
\planteamiento \cref{prop:exp-rango} y \cref{prop:lse}(i).
\pasos
\begin{enumerate}[label=\arabic*.]
  \item (a) $100>88{,}72$, luego $e^{100}=2{,}7\cdot 10^{43}$ desborda a \Inf, igual que $e^{101}$ y $e^{103}$. Cada componente es $\Inf/\Inf=\NaN$.
  \item (b) $\softmax(v-c\uno)_i=\frac{e^{v_i}e^{-c}}{e^{-c}\sum_je^{v_j}}=\softmax(v)_i$, porque $e^{-c}$ sale de la suma y se cancela. $\blacksquare$
  \item (c) $c=103$; $v-c\uno=(-3,-2,0)$; $e^{-3}=0{,}04979$, $e^{-2}=0{,}1353$, $e^{0}=1$; suma $1{,}185$.
  \item $\softmax(v)=(0{,}04979/1{,}185;\ 0{,}1353/1{,}185;\ 1/1{,}185)=(0{,}04201;\,0{,}1142;\,0{,}8438)$.
\end{enumerate}
\resultado \textbf{ingenua: $(\NaN,\NaN,\NaN)$; estable: $(0{,}0420;\,0{,}1142;\,0{,}8438)$.}
\comprobacion las componentes suman $1{,}000$; \texttt{torch.softmax} en fp32 da lo mismo.
\end{solucion}

\begin{solucion}{eje:6-rango}
\datos $x_{\max}$ y $x_{\min}$ de la \cref{tab:formatos}.
\planteamiento \cref{prop:exp-rango}: $t\in[\ln x_{\min},\ln x_{\max}]$.
\pasos
\begin{enumerate}[label=\arabic*.]
  \item fp32: $\ln(3{,}403\cdot 10^{38})=88{,}72$; $\ln 2^{-126}=-126\ln 2=-87{,}34$.
  \item fp16: $\ln 65\,504=11{,}09$; $\ln 2^{-14}=-14\ln 2=-9{,}704$.
  \item bf16: $\ln(3{,}390\cdot 10^{38})=88{,}72$; $\ln 2^{-126}=-87{,}34$.
\end{enumerate}
\resultado \textbf{fp32 y bf16: $t\in[-87{,}34;\ 88{,}72]$; fp16: $t\in[-9{,}70;\ 11{,}09]$. fp16 es el más restrictivo: un logit de~12 ya desborda.}
\comprobacion en PyTorch, $e^{11{,}0}$ en fp16 da $59\,872$ y $e^{11{,}1}$ da \Inf; en fp32, $e^{88{,}7}$ es finito y $e^{88{,}8}$ es \Inf.
\end{solucion}

\begin{solucion}{eje:6-ce}
\datos (a) $z=(2,1,0)$, $y=3$; (b) $z=(0,-20)$ en fp16.
\planteamiento \cref{prop:ce}: $\ell=\LSE(z)-z_y$ y $\nabla_z\ell=\softmax(z)-e_y$.
\pasos
\begin{enumerate}[label=\arabic*.]
  \item (a) $\LSE(z)=2+\ln(1+e^{-1}+e^{-2})=2+0{,}4076=2{,}408$; $\ell=2{,}408-0=2{,}408$.
  \item $\softmax(z)=(0{,}6652;\,0{,}2447;\,0{,}0900)$; $\nabla_z\ell=(0{,}6652;\,0{,}2447;\,-0{,}9100)$.
  \item (b) Ingenua: $e^{-20}=2{,}06\cdot 10^{-9}$ es menor que la mitad del menor subnormal de fp16, $2^{-25}=2{,}98\cdot 10^{-8}$, y se redondea a~0; $\softmax(z)_2=0$ y $\ln 0=-\infty$.
  \item Estable: $\LSE(z)=0+\ln(1+e^{-20})=2{,}06\cdot 10^{-9}$, y $\ln\softmax(z)_2=-20-2{,}06\cdot 10^{-9}=-20{,}000$.
\end{enumerate}
\resultado \textbf{(a) $\ell=2{,}408$ y $\nabla_z\ell=(0{,}665;\,0{,}245;\,-0{,}910)$; (b) ingenua $-\infty$, estable $-20{,}000$.}
\comprobacion las componentes del gradiente suman $0$ (salvo redondeo), porque $\softmax(z)$ y $e_y$ suman~1. PyTorch da \texttt{-inf} con \texttt{log(softmax)} y $-20$ con \texttt{log\_softmax} en fp16.
\end{solucion}

\begin{solucion}{eje:6-cotas}
\datos $z\in\R^{n}$, $c=\max_iz_i$, $S=\sum_ie^{z_i-c}$.
\planteamiento $\LSE(z)=c+\ln S$ (\cref{prop:lse-estable}) y monotonía del logaritmo.
\pasos
\begin{enumerate}[label=\arabic*.]
  \item Cada $z_i-c\le 0$, luego cada término $e^{z_i-c}\in(0,1]$.
  \item Cota inferior: $S$ contiene el término $e^{0}=1$ del máximo y los demás son positivos, así que $S\ge 1$ y $\ln S\ge 0$: $\LSE(z)\ge c$.
  \item Cota superior: $S$ es suma de $n$ términos menores o iguales que~1, así que $S\le n$ y $\LSE(z)\le c+\ln n$.
  \item Igualdad superior: $S=n$ exige que todos los términos valgan~1, es decir, $z_i=c$ para todo $i$; recíprocamente, si todas las componentes son iguales, $S=n$. $\blacksquare$
  \item $z=(3,3)$: $\LSE=3+\ln 2=3{,}693$, la cota superior. $z=(2,1,0)$: $\LSE=2{,}408\in[2;\ 2+\ln 3]=[2;\ 3{,}099]$.
\end{enumerate}
\resultado \textbf{$\max_iz_i\le\LSE(z)\le\max_iz_i+\ln n$, con igualdad superior si y solo si todas las $z_i$ son iguales.}
\comprobacion la cota inferior se alcanza solo en el límite: $\LSE(t,0)\to t$ cuando $t\to+\infty$, porque $\ln(1+e^{-t})\to 0$.
\end{solucion}

\begin{solucion}{eje:6-convexa}
\datos $p=\softmax(z)$; $\nabla\LSE=p$ (\cref{prop:lse}(iii)).
\planteamiento derivar $p_k=e^{z_k}/\sum_ie^{z_i}$ y estudiar la forma cuadrática del hessiano.
\pasos
\begin{enumerate}[label=\arabic*.]
  \item (a) Por la regla del cociente, $\partial p_k/\partial z_j=p_k\delta_{jk}-p_kp_j$. Como $\partial^{2}\LSE/\partial z_j\partial z_k=\partial p_k/\partial z_j$, el hessiano es $\diag(p)-pp\T$.
  \item (b) $v\T\diag(p)v=\sum_ip_iv_i^{2}$ y $v\T pp\T v=(\sum_ip_iv_i)^{2}=\bar v^{2}$.
  \item Usando $\sum_ip_i=1$: $\sum_ip_i(v_i-\bar v)^{2}=\sum_ip_iv_i^{2}-2\bar v\sum_ip_iv_i+\bar v^{2}=\sum_ip_iv_i^{2}-\bar v^{2}$, que es la forma cuadrática. Es una suma de términos no negativos.
  \item Un hessiano semidefinido positivo en todo punto implica convexidad (módulo~4). La entropía cruzada $\ell=\LSE(z)-z_y$ es la suma de una función convexa y una lineal, luego es convexa.
  \item (c) Con $v=\uno$, $\bar v=1$ y $v_i-\bar v=0$ para todo $i$: la forma cuadrática es nula y $(\diag(p)-pp\T)\uno=p-p(p\T\uno)=0$. Es coherente con la \cref{prop:lse}(i): a lo largo de la dirección $\uno$, $\LSE$ es lineal y su segunda derivada se anula.
\end{enumerate}
\resultado \textbf{$\nabla^{2}\LSE=\diag(p)-pp\T\succeq 0$, con $\uno$ en su núcleo; $\LSE$ y la entropía cruzada son convexas en los logits.} $\blacksquare$
\comprobacion para $n=2$ y $z=(0,0)$, $p=(0{,}5;\,0{,}5)$ y el hessiano es $\begin{psmallmatrix}0{,}25&-0{,}25\\-0{,}25&0{,}25\end{psmallmatrix}$, con autovalores $0{,}5$ y $0$.
\end{solucion}
